\id{ҒТАМР 65.53.33}{}

\begin{header}
\swa{}{БАҚША ЖӘНЕ КӨКӨНІС ДАҚЫЛДАРЫН КЕПТІРУДІҢ ОҢТАЙЛЫ РЕЖИМДЕРІН ЖӘНЕ АЛЫНҒАН ҰНТАҚТАРДЫҢ СИПАТТАМАЛАРЫН ЗЕРТТЕУ}

Г.Е.~Жумалиева\envelope,
У.Ч.~Чоманов,
А.Б.~Байзақова,
А.К.~Шоман
\end{header}

\begin{affil}
«Қазақ қайта өңдеу және тағам өнеркәсіптері ғылыми-зерттеу институты» ЖШС Алматы қаласы филиалы, Казақстан

\corrauthor{Корреспонден-тавтор: guljan\_7171@mail.ru}
\end{affil}

Бұл зерттеуде лиофильді және инфрақызыл кептіру арқылы алынған бақша
және көкөніс дақылдары ұнтақтарын кешенді зерттеу жүргізілді. Зерттеудің
мақсаты --- дайын өнімдердің технологиялық, органолептикалық,
физико-химиялық, түс және витаминдік қасиеттеріне әртүрлі кептіру
режимдерінің әсерін бағалау болды.

Зерттеу нысандары ретінде отандық қызанақ пен қарбыз сорттары алынды.
Ұнтақтар әртүрлі температуралық және вакуумдық режимдерде лиофильді
кептіру және инфрақызыл кептіру арқылы алынды. Алынған үлгілердің
ылғалдылығы, титрленетін және активті қышқылдығы, органолептикалық
қасиеттері, CIE Lab* жүйесі бойынша түс сипаттамалары және жоғары өнімді
сұйық хроматография әдісі арқылы витаминдік құрамы зерттелді.

Зерттеу нәтижесінде өнімнің құрылымы, табиғи түсі және дәмі сақталған,
қалдық ылғалдылығы 12,5--13,4 \% болатын ұнтақтар алынды. Лиофильді
кептіру режимі --40\,°C және --55\,°C температурада, ал инфрақызыл
кептіру қызанақ үшін 60\,°C, қарбыз үшін 50\,°C температурада
жүргізілгенде ұнтақтар жоғары түстік қанықтылыққа және түс тұрақтылығына
ие болды. Витаминдік талдау А, Е, С және В дәрумендерінің жоғары
деңгейін көрсетті, бұл өнімнің биологиялық құндылығын растайды.

Зерттеу нәтижелері мұздатып кептіру және инфрақызыл кептіру әдістерін
қолданудың қауын мен қызанақ сияқты бақша дақылдарын өңдеуде және
функционалды ұнтақ өнімдерін алуда перспективалы екенін көрсетеді.

{\bfseries Түйін сөздер:} қарбыз, қызанақ, ұнтақтар, бақша мен көкөніс
дақылдары, инфрақызыл кептіру, лиофильді кептіру.

\begin{header}
ИССЛЕДОВАНИЕ ОПТИМАЛЬНЫХ РЕЖИМОВ СУШКИ БАХЧЕВЫХ И ОВОЩНЫХ КУЛЬТУР И ХАРАКТЕРИСТИК ПОЛУЧЕННЫХ ПОРОШКОВ

Г.Е.~Жумалиева\envelope,
У.Ч.~Чоманов,
А.Б.~Байзакова,
А.К.~Шоман
\end{header}

\begin{affil}
Филиал ТОО «Казахский научно-исследовательский институт перерабатывающей и пищевой промышленности» в г. Алматы, Казахстан,

e-mail: guljan\_7171@mail.ru
\end{affil}

В данном исследовании проведено комплексное изучение порошков из
бахчевых и овощных культур, полученных методом лиофильной и инфракрасной
сушки. Целью исследования являлась оценка влияния различных режимов
сушки на технологические, органолептические, физико-хими\-ческие, цветовые
и витаминные свойства готовой продукции.

В качестве объектов исследования были выбраны отечественные сорта томата
и арбуза. Порошки получали при различных температурных и вакуумных
режимах лиофильной и инфракрасной сушки. Полученные образцы
анализировали по показателям влажности, титруемой и активной
кислотности, органолептическим свойствам, цветовым характеристикам по
системе CIE Lab*, а также по витаминному составу с использованием метода
высокоэффективной жидкостной хроматографии.

В результате исследования были получены порошки с остаточной влажностью
12,5--13,4 \%, характеризующиеся сохранением структуры, естественного
цвета и вкуса продукта. Установлено, что при лиофильной сушке при
температурах --40 °C и --55 °C, а также при инфракрасной сушке при
температуре 60 °C для томата и 50 °C для арбуза, порошки обладают
высокой цветовой насыщенностью и стабильностью окраски. Анализ
витаминного состава показал высокое содержание витаминов A, E, C и
группы B, что подтверждает высокую биологическую ценность продукции.

Полученные результаты свидетельствуют о перспективности применения
лиофильной и инфракрасной сушки для переработки бахчевых и овощных
культур, таких как арбуз и томат, и получения функциональных
порошкообразных продуктов.

{\bfseries Ключевые слова:} арбуз, томат, порошки, бахчевые и овощные
культуры, инфракрасная сушка, лиофильная сушка.

\begin{header}
STUDY OF OPTIMAL DRYING REGIMES FOR MELON AND VEGETABLE CROPS AND THE CHARACTERISTICS OF THE OBTAINED POWDERS

G.~Zhumalieva\envelope,
U.~Chomanov,
A.~Baizakova,
A.~Shoman
\end{header}

\begin{affil}
Branch of LLP ``Kazakh Research Institute of Processing and Food Industry'' in Almaty,

e-mail: guljan\_7171@mail.ru
\end{affil}

In this study, a comprehensive investigation of powders obtained from
melon and vegetable crops using freeze-drying and infrared drying was
conducted. The aim of the study was to evaluate the effect of different
drying regimes on the technological, organoleptic, physicochemical,
color, and vitamin properties of the final products.

Domestic tomato and watermelon varieties were selected as the research
objects. The powders were produced under various temperature and vacuum
conditions using freeze-drying and infrared drying methods. The obtained
samples were analyzed for moisture content, titratable and active
acidity, organo\-leptic properties, color characteristics using the CIE
Lab* system, and vitamin composition by high-performance liquid
chromatography.

The results showed that powders with a residual moisture content of
12.5--13.4\% were obtained while preserving the structure, natural
color, and taste of the products. Freeze-drying at --40 °C and --55 °C
and infrared drying at 60 °C for tomato and 50 °C for watermelon
provided powders with high color saturation and color stability. Vitamin
analysis revealed high levels of vitamins A, E, C, and B-group vitamins,
confirming the high biological value of the products.

The results demonstrate the potential of freeze-drying and infrared
drying for processing melon and vegetable crops, such as watermelon and
tomato, and for producing functional powdered products.

{\bfseries Keywords:} watermelon, tomato, powders, melon and vegetable
crops, infrared drying, freeze-drying.

\begin{multicols}{2}
{\bfseries Кіріспе}. Тамақтану - ағзаның қалыпты жұмыс істеуі мен
өнімділігін және денсаулықы сақтаудың негізгі факторы болып табылады.
Ғылымда 19 ғасырдың аяғы мен 20 ғасырдың басында жүргізілген зерттеулер,
біздің қазіргі заманғы қоректік қажеттіліктер туралы түсінігімізді
қалыптастырды. Барлығына белгілі маңызды қоректік заттар ретінде -
аминқышқылдары, май қышқылдары, дәрумендер және минералдар болып
саналады.20 ғасырдың басында негізі қаланған тамақтану теориясы бүгінгі
күнге дейін өзінің өзектілігін сақтаған. Бұл теорияның негізгі қағидасы:
тамақтану - бұл ағзаның химиялық және биологиялық құрамын сақтау және
теңестіру процесі. Уақытылы және сапалы тамақтану ағзаның қалыпты жұмыс
істеуін қамтамасыз ететін тағам компоненттерінің оңтайлы арақатынасын
білдіреді.

Ағзаның бейімделу қасиетін сақтау үшін тағам рационында ақуыздар,
дәрумендер және биологиялық белсенді қосылыстар және макро- және
микроэлементтер болуы керек.

Қазіргі уақытта медициналық зерттеулер қатерлі ісікке қарсы дәрі
іздеуде. Химиялық препараттардан басқа, зерттеушілер табиғи антиоксидант
ликопинді қатерлі ісік қаупінің төмендеуіне ықпалын тигізетіні анықтады.
2001 жылы жүргізілген клиникалық сынақтың алғашқы нәтижелері ликопиннің
қуық асты безі қатерлі ісігі жасушаларының өсу қарқынын төмендететінін
көрсетті. Бұл әдіс қатерлі жасушалардың таралуын 73\%-ға дейін азайтқан.
З ерттеулерде 46 түрлі жеміс-жидек пен көкөністің ішінде тек қызанақ
қатерлі ісік қаупінің төмендеуімен айқын әсер ететіні байқалған {[}1{]}.

Адам тағам құрамында кездесетін 60 каротиноидтың ішіндегі ең маңыздысы-
ликопин және бета-каротин. Бұл каротиноидтардың ерекшелігі, олардың
хромофоры, ол қос байланыстардың ұзын тізбегінен тұрады, бұл оларды
табиғи пигменттерге айналдырады. Кейбір өсімдіктерде, басқа
каротоидтарнмен қатар, белгілі бір каротиноидтың басым концентрациясы
бар (мысалы, қызанақта ликопин, сары жүгерідегі криптоксантин және
т.б.).

Ликопин (C₄₀H₅₆) - қызыл жемістер мен көкөністерде (қызанақ, бұрыш,
қарбыз, қызғылт грейпфрут, өрік) кездесетін каротиноидтар тобына жататын
пигмент, бұл өсімдік текті тағамдарға қызыл түс береді. Суда ерімейді
және майларда ериді. Ликопиннің антиоксиданттық белсенділігі
бета-каротиндерден үш есе жоғары {[}2{]}. Піспенан өндірісінде
функционалды ингредиент ретінде қызанақ және қарбыз ұнтағын таңдау келсі
факторларға негізделген: республикадағы ауыл шаруашылығы және қайта
өңдеу кәсіпорындарында өсімдік материалдары мен шырын өндірісінің
қалдықтарынан тамақ өнімдеріне арналған ұнтақтар өндіру, қалдықсыз
өндірісті дамытуға мүмкіндік береді. Қызанақты және қарбызды өңдеу
кезінде пайда болатын қосалқы өнімдер тағамдық қоспаларды өндіру үшін
құнды шикізат болып табылады {[}3{]}.

Қызанақтың химиялық құрамы жағынан ерекше және сортына, өсіру
жағдайларына және технологиясына байланысты өзгеріске аз ұшырайды.
Құрамындағы оңай сіңетін макронутриенттер, орташа есеппен 16,5 г
тағамдық талшық және барлық маңызды аминқышқылдары бар ақуыздар және
омега-3 және омега-6 май қышқылдарына бай келеді. Калий, кальций,
магний, натрий, фосфор, темір, марганец, мыс, селен және мырышқа бай.
Витаминдерден қызанақ құрамында В дәрумендері, С, А, Е және К
дәрумендері, сондай-ақ каротиноидтар бар, олардың басым бөлігі ликопин
болып табылады {[}4, 5{]}.

Ликопин қызанақ ұнтағына қызыл түс беріп қана қоймай, сонымен қатар
күшті табиғи антиоксидант. Тағамдық қышқылдардың орташа мөлшері 0,5\%,
тағамдық талшық 0,8\% және пектин 0,3\% құрайды. Қызанақтар Е дәрумені,
ниацин және пантотен қышқылының ең көп мөлшері бар көкөністердің
қатарына жатады, ал пектин холестеринді төмендетуге көмектеседі{[}6{]}.

Қызанақ антидепрессанттық әсерге ие, серотониннің болуына байланысты
жүйке жүйесін реттейді, сондай-ақ бактерияға қарсы және қабынуға қарсы
қасиеттерге ие. Осы қасиеттердің барлығын сақтай отырып, қызанақ ұнтағы
табиғи бояғыш және хош иістендіргіш ретінде кеңінен қолданылады. Ол
өнімдерге қышқыл дәм, айқын хош иіс және қанық қызыл түс береді{[}7{]}.

Қызанақ және қарбыз ұнтақтары, сондай-ақ олардан жасалған қоспалар
ликопинге, цитруллинге, дәрумендерге және минералдарға бай, бұл оларды
адам денсаулығын сақтаудың тиімді құралы етеді. Олар антиоксидантты,
қабынуға қарсы және кардиопротекторлық қасиеттерге ие, жүрек-қан
тамырлары аурулары мен қатерлі ісіктің алдын алуға және ылғалдануды
сақтауға көмектеседі.

Қарбыз және қызанақ ұнтағы - ликопинге, дәрумендер мен минералдардың
қоймасы болып табылатын концентрлі өнім. Ликопин бос радикалдарды
бейтараптандырады, жасушалардың зақымдануын тежейді және қабынуды
азайтады. Ол жүрек-қан тамырлары ауруларының, атеросклероздың және
қатерлі ісіктің алдын алуда маңызды рөл атқарады{[}8{]}.

Қарбыз ұнтағы ликопинмен қатар, С және А дәрумендеріне және цитруллинге
де бай. Цитруллин қан айналымын жақсартады және қан қысымын төмендетеді,
бұл қарбыз ұнтағын гипертониясы бар адамдар үшін пайдалы етеді. Ол
сондай-ақ су-тұз балансын реттейді және гидратацияны жақсартады.

{\bfseries Материалдар мен әдістер}. Зерттеу нысандары ретінде қарбыз бен
қызанақ алынды.

Зерттеу барысында стандарттарға сәйкес келесі әдістемелер қолданылды:

- МЕМСТ ISO 750-2013 - Жеміс және көкөніс өнімдерін өңдеу
өнімдері;

- МЕМСТ 29031-91 - Құрғақ заттардың мөлшерін анықтау әдісі.
Құрғақ заттардың мөлшері СНЕЛ-104 рефрактометрінде анықталды. Шикізат
пен дайын өнімнің ылғалдылығы МХ-50 ылғалөлшегішімен өлшенді. Өлшеу
нұсқаулыққа сәйкес жүргізілді: 5 г сынама алынып, табақшаға біркелкі
жайылып, 130 °C температурада 40 минут кептірілді;

- МЕМСТ 34130-2017 - Қызанақ ұнтағының
органолептикалық сапа көрсеткіштерін анықтау;

- МЕМСТ EN 12823-2-2014 - Тағам өнімдеріндегі ликопин мөлшерін
жоғары өнімді сұйық хроматография әдісімен анықтау;

- МЕМСТ 5898--87 - Қышқылдық пен сілтілікті анықтау;

- МЕМСТ Р 51478-99 (ISO 2917-74) - Сутек иондары
концентрациясын (pH) анықтау;

- МЕМСТ Р ISO 22935-2-2011 - Органолептикалық бағалау әдістері;

- testo 206-pH1 pH-метрі арқылы активті қышқылдықты анықтау;

- МЕМСТ EN 12823-2-2014 - А витаминінің мөлшерін жоғары өнімді
сұйық хроматография әдісімен анықтау;

- МЕМСТ EN 12822-2014 - \emph{Е витаминінің (α-, β-, γ-,
δ-токоферолдар) мөлшерін жоғары өнімді сұйық хроматография әдісімен
анықтау;}

- МЕМСТ 34151-2017 - \emph{С витаминін жоғары өнімді
сұйық хроматография әдісімен анықтау;}

- бақша және көкөніс дақылдарынан кептірілген ұнтақтарды алу әдістері.

Ликопинге бай ұнтақтарды бақша және көкөніс дақылдарынан алу үшін екі
әдіс қарастырылды, лиофильді кептіру және инфрақызыл кептіру. Қызанақ
пен қарбыз массасынан артық ылғалды жою үшін TAGLER ЦЛМН 1-8 зертханалық
центрифугасы қолданылды.

\emph{Лиофильді кептіру} - өнімдегі суды төмен температура мен төмен
қысымда сублимация арқылы (мұздың тікелей буға айналуы) жою процесі. Бұл
әдіс ликопин сияқты қоректік заттар мен биологиялық белсенді
қосылыстардың көп бөлігін сақтауға мүмкіндік береді, себебі процесс
төмен температурада өтеді, бұл сезімтал компоненттердің термиялық
бұзылуын азайтады. Лиофильді кептіру сондай-ақ бастапқы өнімнің
құрылымын және дәмін сақтайды, сонымен қатар оның массасын едәуір
азайтады, бұл тасымалдау және сақтау процесін жеңілдетеді. Бұл процесс
бірнеше негізгі кезеңнен тұрады:

- алдын ала мұздату - шикізат төмен температураға дейін мұздатылады, бұл
жасуша ішінде мұз кристалдарының түзілуіне ықпал етеді.

- біріншілік (сублимациялық) кептіру - төмен қысымда бос ылғалды жою,
мұнда су қатты күйден (мұз) бу күйіне өтеді, сұйық фазадан өтпей.

- екіншілік (десорбциялық) кептіру - байланысқан ылғалды одан әрі төмен
қысым мен температурада жою.

Лиофильді кептіру биологиялық және фармацевтикалық препараттарды,
сондай-ақ тағам өнеркәсібінде жемістер, көкөністер және басқа өнімдерді
кептіру үшін белсенді қолданылады.

Зертханалық жағдайда лиофилизация процесін Martin Christ Alpha 1-2 LD
Plus лиофильді кептіргішімен жүзеге асырылды.

\emph{Инфрақызыл кептіру} инфрақызыл сәулеленуді қолдану арқылы
материалдан ылғалды жоюға негізделген. Инфрақызыл сәулелену су
молекулаларына тікелей әсер етіп, олардың қызуынан ылғалдың тез
булануына әкеледі. Бұл әдіс дәстүрлі кептіру әдістеріне қарағанда
жылдамырақ және энергияны аз қажет етеді, мысалы, ыстық ауада кептіру
немесе тіпті лиофилизация. Сонымен қатар, бұл әдіс органолептикалық
қасиеттерді, мысалы, түс пен дәмді сақтауға ықпал етеді және жақсы
дәмдік қасиеттері бар, салыстырмалы түрде жоғары ликопин концентрациясы
бар ұнтақтарды тиімді алуға мүмкіндік береді.

Жемістер үшін Basic Station 3 инфрақызыл кептіргіш шкафы пайдаланылды.

Қызанақ және қарбыз ұнтақтары ірі жемісті қызанақтарды ұсақтау, IKA A11
basic (ИКА, Германия) аналитикалық ұнтақтағыш арқылы алынды.

{\bfseries Нәтижелер және талқылау}. Зерттеуде сублимациялық кептіру және
инфрақызыл кептіру арқылы жаңа піскен қызанақтар мен қарбыздардан
ликопинге бай құрғақ ұнтақтарды алу әдістері қолданылды. Бұл әдіс
ликопин сияқты биологиялық белсенді заттарды сақтау тиімділігін
қамтамасыз етеді, сондай-ақ әртүрлі технологиялардың осы заттардың
сақталуына және алынған ұнтақтардағы ликопин мөлшеріне әсерін зерттеуге
мүмкіндік береді.

Зерттеу нысандары ретінде қызанақ пен қарбыздың отандық сорттары
пайдаланылды.

Ликопинге бай қызанақ пен қарбыздан құрғақ ұнтақтарды алудың екі әдісі
қарастырылды: жеміс-жидек және көкөніс ұнтақтарын алу үшін тамақ
өнеркәсібінде кеңінен қолданылатын мұздатып кептіру (лиофильді) және
инфрақызыл кептіру.

Лиофилизация төмен температурада және төмен қысымда жүзеге асырылады,
бұл мұзда сублимация арқылы суды кетіреді, жылу әсерін азайтады және
осылайша ликопинді қоса алғанда, термосезімтал заттардың максималды
мөлшерін сақтайды. Инфрақызыл кептіру, керісінше, материалды инфрақызыл
сәулелену арқылы қыздырады, бұл кептіру процесін жеделдетеді, бірақ жылу
әсерінен биологиялық белсенді заттардың жоғалуы жоғары болуы мүмкін.

Ликопинді сақтау және алынған ұнтақтардың сапасын бағалау үшін жоғарыда
аталған әдістердің тиімділігі анықталды және жоғары тағамдық және
функционалдық құндылығы бар ликопинге бай өнімдерді алудың ең тиімді
әдісі таңдалды.

\emph{Қызанақ пен қарбыздан құрғақ ұнтақты лиофилизация әдісімен алудың
технологиялық процесі}

Шикізат дайындау үшін қызанақ пен қарбызды, артық ылғалды кетіру үшін
центрифугаланды. Шикізат ретінде жаңа піскен қызанақ пен қарбыздар
алынып, мұқият жуылып, кептірілді. Содан кейін жемістер біртекті масса
алу үшін блендерде араластырылды. Алынған массаны центрифугаға салып,
25°C температурада және 15 000 айн/мин жылдамдықта 20 минут
центрифугаланды.

Центрифугалау нәтижесінде бастапқы ликопин мен басқа да биологиялық
белсенді қосылыстары бар қызанақ және қарбыз целлюлозасынан тұратын
концентрлі тұнба алынды. Артық сұйықтық тиімді түрде жойылды. Алынған
қызанақ пен қарбыз целлюлозасы 10-25 мм қалыңдықтағы қабатта арнайы
табаларға біркелкі орналастырылып, лиофилизацияға жіберілді

Алдын ала мұздату -10°C температурада 12 сағат бойы жүргізілді, себебі
бұл кезең лиофилизацияға арналған үлгілерді дайындау үшін қажет және
материалдан суды тиімді түрде кетіру үшін мұз кристалдарының пайда
болуын қамтамасыз етеді.

1. Тәжірибе дайын өнімнің қалдық ылғалдылығын 10-12\% деңгейінде ұстау
мақсатында үш нұсқада жүргізілді. Кептіру ұзақтығы әр нұсқаның жағдайына
байланысты 24-тен 48 сағатқа дейін өзгеріп отырылды, бұл кезеңде өнімнің
ылғалдылығы біртіндеп төмендеді.

Қызанақ пен қарбыз үшін үш түрлі кептіру режимі таңдалды, олардың
әрқайсысы температура мен вакуумдық әсерді екі кезеңде біріктіреді. Бұл
режимдер зерттеліп жатқан өнімдердің суды ұстау қабілеті мен жылуға
сезімталдығын ескере отырып, әртүрлі параметрлердің ылғалды кетіру
тиімділігіне және биологиялық белсенді заттардың сақталуына әсерін
бағалау үшін әзірленді. Режимдер теориялық сублимациялық кептіру
болжамдары және төмен температура мен қысымда өсімдік тіндерінің
мінез-құлқын зерттеу нәтижелеріне негізделген.

{\bfseries Нұсқа 1:} кезең 1 - -- 30°C, 0,37 mBar; кезең 2 --- --45°C, 0,1
mBar.

{\bfseries Нұсқа 2:} кезең 1 - -- 35°C, 0,5 mBar, 10 сағ; кезең 2 ---
--55°C, 0,2 mBar.

{\bfseries Нұсқа 3:} кезең 1 - -- 40°C, 0,3 mBar; кезең 2 --- --55°C, 0,05
mBar.

Бірінші нұсқадағы кептіру режимі (температура -30°C, қысым 0,37 мБар,
ұзақтығы 12 сағат) бірінші кезеңде еріген суды мұз кристалдары түрінде
мұздатуды қамтамасыз етті және оны сублимация арқылы тиімді түрде
кетіруге мүмкіндік берді. Екінші кезеңде температура -45°C дейін, қысым
0,1 мБар, ұзақтығы 10 сағатқа дейін төмендетілді, бұл өнім құрылымында
берік ұсталатын байланысқан ылғалға әсер етуге мүмкіндік берді.

Екінші нұсқада бірінші кезеңде 10 сағат бойы -35°C төмен температура,
бірақ 0,5 мБар жоғары қысым қолданылды. Бұл жағдайлар бос судың
кристалдануы мен сублимациясы үшін онша тиімді болмады, себебі
мұздатудың жеткіліксіздігі және жоғары қысым фазалық ауысу процесін
баяулатуы мүмкін. Екінші кезеңде -55°C температура, қысым 0,2 мБар, 10
сағат - қосымша ылғалды кетіруді қамтамасыз етті.

Үшінші нұсқа бірінші кезеңде температураны (-40°C) төмендету және
қысымды (0,3 мБар) төмендету арқылы мұздату және сублимация процесін
жеделдетуге бағытталған, бірақ кезеңнің ұзақтығы 8 сағатқа дейін
қысқартылды. Жылдам мұздату кептіру жылдамдығына оң әсер етуі мүмкін,
бірақ төмен температурамен жеткілікті уақыт болмаған жағдайда бос су
толығымен жойылмауы мүмкін. Екінші кезеңде -55°C температура, 0,05 мБар
қысым, 10 сағат - байланысқан ылғалды тиімді түрде жояды, бірақ
қарбыздағы қант пен органикалық қышқылдардың көптігіне байланысты кейбір
сезімтал компоненттердің ыдырауына әкелді.

Қызанақтағы ылғалдың өзгеруінің нәтижелері 1-кестеде, ал қарбызда
2-кестеде көрсетілген.
\end{multicols}

\tcap{1-кесте. Лиофильді кептіру процесінде қызанақ ұнтағының ылғалдылығының уақытқа байланысты өзгеруі}
\begin{longtblr}[
  label = none,
  entry = none,
]{
  cells = {c},
  cells = {font = \small},
  row{1} = {font=\bfseries\small},
  cell{1}{1} = {r=3}{},
  cell{1}{2} = {c=3}{0.514\linewidth},
  cell{2}{2} = {c=3}{0.514\linewidth},
  hlines,
  vlines,
}
Уақыт ұзақтығы, сағат & Лиофильді кептіру кезінде қызанақ ұнтағының ылғалдылығының өзгеруі &  & \\
 & Ылғалдылығы, W, \% &  & \\
 & нұсқа 1 & нұсқа 2 & нұсқа 3\\
12 & 55,3 & 45,5 & 35,9\\
14 & 30,7 & 36,5 & 28,6\\
16 & 25,4 & 27,7 & 22,5\\
18 & 17,1 & 20,6 & 13,3\\
20 & 14,9 & 11,5 & 7,41
\end{longtblr}

\begin{multicols}{2}
1-кестеде келтірілген нәтижелерге сәйкес, ең жоғары тиімділік 3-нұсқада
байқалды, себебі 12 сағат кептіргеннен кейін өнімнің ылғалдылығы 35,9\%
құрады, 1-нұсқада 55,3\% және 2-нұсқада 45,5\% болды.18-ші сағатқа
қарай ұнтақтың ылғалдылығы 3-нұсқада 13,3\% болды, 1 және 2-нұсқаларда
сәйкесінше 17,1\% және 20,6\% болды. Максималды сусыздану кептірудің
20-шы сағатында байқалды: ең төменгі қалдық ылғалдылық 3-нұсқада 7,41\%
болды, бұл судың ең толық жойылуын көрсетеді.

2-нұсқадағы кептіру режимі қызанақ шикізатынан бос және байланысқан
ылғалды кетіру үшін температура мен вакуум жағдайларының ең тиімді
үйлесімін көрсетті.100 г қызанақ сығындысынан алынған құрғақ зат
мөлшері 4,966 г құрады, бұл қызанақ жемістеріндегі құрғақ заттың
әдеттегі мөлшеріне сәйкес келеді.
\end{multicols}

\tcap{2-кесте. Лиофильді кептіру процесінде қарбыз ұнтағының ылғалдылығының уақытқа байланысты өзгеруі}
\begin{longtblr}[
  label = none,
  entry = none,
]{
  cells = {c},
  cells = {font = \small},
  row{1} = {font=\bfseries\small},
  cell{1}{1} = {r=3}{},
  cell{1}{2} = {c=3}{0.514\linewidth},
  cell{2}{2} = {c=3}{0.514\linewidth},
  hlines,
  vlines,
}
Уақыт ұзақтығы, сағат & Лиофильді кептіру кезінде қарбыз ұнтағының ылғалдылығының өзгеруі &  & \\
 & Ылғалдылығы,W, \% &  & \\
 & нұсқа 1 & нұсқа 2 & нұсқа 3\\
27 & 46,8 & 54,2 & 65,7\\
29 & 40,6 & 41,8 & 43,9\\
31 & 35,2 & 33,5 & 24,1\\
33 & 29,8 & 21,7 & 12,5\\
35 & 20,4 & 13,7 & 9,3\\
37 & 14,2 & 10,4 & 4,7
\end{longtblr}

\begin{multicols}{2}
Қарбыз ұнтағы үшін (2-кесте) ылғалды азайту динамикасы тұрақты төмен
қалдық ылғалдылыққа қол жеткізу үшін көбірек уақыт қажет екенін
көрсетеді, бұл негізінен қарбыз целлюлозасының су мөлшерінің жоғары
болуына байланысты.3-нұсқа ең жоғары кептіру тиімділігін көрсетті: 33
сағат кептіргеннен кейін ылғалдылық 12,5\%-ға дейін төмендеді, ал 37
сағаттан кейін 4,7\%-ға жетті.

1- нұсқада 33 және 37 сағаттан кейінгі ылғалдылық сәйкесінше 29,8\% және
14,2\% болды, ал 2-нұсқада 21,7\% және 10,4\% болды. Бұл режимде 3-нұсқа
(--40°C және --55°C, қысымның төмендеуі және кептіру уақыты жеткілікті)
қарбыз ұнтағын тиімді түрде құрғатты.100 г қарбыз сығындысынан алынған
құрғақ заттың мөлшері 3,132 г болды, бұл қызанақпен салыстырғанда
қарбыздағы құрғақ заттың аз мөлшеріне сәйкес келеді.

Нәтижелер көрсеткендей, -40°C және -55°C температурада сатылы кептірумен
3-нұсқа қызанақ пен қарбыз үшін ең тиімді кептіру әдісі болды. Бұл әдіс
бойынша қызанақ ұнтағының ылғалдылығы 18 сағат кептіргеннен кейін
13,3\%-ға төмендеді, ал қарбыз ұнтағы үшін 33 сағат кептіргеннен кейін
ылғалдылық 12,5\%-ға жетті.

Мұздатып кептіру өнімнің құрылымын және органолептикалық қасиеттерін,
соның ішінде түсін, дәмін және хош иісін сақтауға көмектесті, осылайша
жоғары сапалы, антиоксиданттарға бай қызанақ ұнтағын өндіруді
оңтайландырды.

\emph{Қызанақ және қарбыз ұнтақтарын инфрақызыл кептіру әдісімен алу
технологиялық процесі.}

Инфрақызыл кептіру арқылы жаңа піскен қызанақтарды алу әдісі, шикізатты
ұсақтау және кептіру арқылы келесідей болды: қызанақтар тексеріліп,
жуылып, сабақтары алынды. Инфрақызыл өңдеу және кептіру кезінде
қызанақтардан ылғалды тиімді кетіру үшін қызанақтар сегіз тең тілімге
кесілді.

Қызанақтарды осылай кесудің себебі, алюминий тағам табағына терісі
жоғары қаратып қойылған кезде, тілімдер кептіру кезінде табаға жабысып
қалмайды, бұл сұйықтықтың шырын камерасынан төменде орналасқан
инфрақызыл қыздырғыштарға ағуына жол бермейді.

Тілімдер бір қабатқа біркелкі орналастырылып, арасы 0,5-1,5 см болды.
Содан кейін таба кіріктірілген керамикалық инфрақызыл қыздырғыштары бар
кептіру пешіне салынды. Қызанақтар инфрақызыл кептіргіште 45°C, 50°C
және 60°C температурада ылғалдылық 10-12\% жеткенше кептірілді.

Инфрақызыл сәулеленуді қолдана отырып, қарбыз ұнтағын алу әдісі, оның
ішінде шикізатты ұнтақтау және кептіру келесідей жүзеге асырылды:
қарбыздар тексеріліп, жуылып, аршылып, тұқымдары алынып, салмағы
0,014--0,022 кг болатын кесектерге кесілді. Содан кейін кесектер жалпақ
пісіру парағына бір қабат етіп салынды.

Пісіру парағы кіріктірілген керамикалық инфрақызыл жылытқыштары бар
кептіру пешіне қойылды. Инфрақызыл энергияны қолдана отырып, үш кептіру
режимі зерттелді: қарбыздарды 40, 50 және 60°C температурада қыздыру.

Камерадағы жұмыс температурасы 40--60°C аралығында орнатылды, ал
қызанақтар осы температурада 18,0--24,5 сағат ұсталды. Содан кейін
камерадағы жұмыс температурасы 20--30°C дейін төмендетілді, ал
қызанақтар осы температурада 2,5--3,0 сағат ұсталды.

Кептіргеннен кейін кептіру камерасынан алынып, қызанақтар ұнтақ алу үшін
аналитикалық ұнтақтағышта ұнтақталды. Қыздырудың ең жоғары рұқсат
етілген температурасы 60°C болып белгіленді, себебі бұл температурадан
жоғары температура кептірілген қызанақтардың өзгеруін айтарлықтай
арттырып, олардың сапасына кері әсер етті. Инфрақызыл кептіру нәтижеле
\end{multicols}

\tcap{3-кесте. Инфрақызыл кептіру режимдері және ұнтақ ылғалдылығының өзгеруі}
\begin{longtblr}[
  label = none,
  entry = none,
]{
  cells = {c},
  cells = {font = \small},
  row{1} = {font=\bfseries\small},
  row{2} = {font=\bfseries\small},
  cell{1}{1} = {r=2}{},
  cell{1}{2} = {c=6}{0.414\linewidth},
  cell{2}{2} = {c=3}{0.212\linewidth},
  cell{2}{5} = {c=3}{0.201\linewidth},
  hlines,
  vlines,
}
Көрсеткіштер & Зерттелген ұнтақтар &  &  &  &  & \\
 & Қызанақ ұнтағы &  &  & Қарбыз ұнтағы &  & \\
Қыздыру температурасы, °С & 40 & 50 & 60 & 40 & 50 & 60\\
Оптималды ылғалдылыққа жету уақыты, час & 23 & 21 & 20 & 39 & 37 & 34\\
Ылғалдылық, \% & 12,4 & 13,6 & 13,1 & 13,9 & 13,6 & 13,4
\end{longtblr}

\begin{multicols}{2}
Тәжірибелік деректерге сүйене отырып, қосымшадағы 3-кестеде
көрсетілгендей, ылғалдылығы 10-13\% болатын қызанақ және қарбыз
ұнтақтарын алу үшін оңтайлы инфрақызыл кептіру температуралары
анықталды. Әртүрлі температурадағы кептіру уақытын талдау, сондай-ақ
дайын өнімнің органолептикалық қасиеттері мен сыртқы түрін зерттеу
арқылы ұтымды процесс параметрлері белгіленді.

Қызанақ ұнтағы үшін кептіру уақыты қыздыру температурасының
жоғарылауымен қысқарды: 40°C температурада кептіру уақыты 23 сағат, 50°C
температурада 21 сағат және 60°C температурада 20 сағат болды.60°C
температураны пайдалану органолептикалық қасиеттердің нашарлауына
әкелген жоқ. Атап айтқанда, ұнтақтың түсі қанық, ашық қызыл болып қалды,
табиғи реңкінің қараңғылануы немесе жоғалуы болмады. Өнімнің дәмі мен
хош иісі тұрақты болды және кептірілген өсімдік материалдарына қойылатын
сапа талаптарына сай болды. Бұл бақылаулар көрсеткендей, қызанақ
шикізатын 60°C температурада кептіру оңтайлы болып саналады, себебі бұл
температура органолептикалық қасиеттерді сақтай отырып және термиялық
зақымдану белгілерін болдырмай, қысқа кептіру уақытын (20 сағат)
қамтамасыз етеді.

Қызанақ ұнтағынан айырмашылығы, қарбыз ұнтағы үшін температураны 60°C-қа
дейін көтеру өнім сапасына кері әсерін тигізді. Кептіру уақыты 40°C
температурада 40 сағаттан 60°C температурада 34 сағатқа дейін
қысқартылғанымен, максималды температурада ұнтақтың айтарлықтай қараюы
байқалды, бұл шикізат компоненттерінің термиялық ыдырауының басталғанын
көрсетеді. Түстің айтарлықтай жоғалуы және дәм мен хош иістің нашарлауы
бұл температураны практикалық емес етеді. Қарбыз ұнтағы үшін ең қолайлы
температура - 50°C, кептіру уақыты - 37 сағат. Бұл жағдайда
органолептикалық қасиеттердің тұрақтылығы қамтамасыз етілді: ұнтақ
өзінің табиғи түсін, хош иісін және дәмін сақтап қалды, бұл термиялық
зақымданудың жоқтығын растайды.

Қызанақ ұнтағынан айырмашылығы, қарбыз ұнтағы үшін температураны 60°C-қа
дейін көтеру өнім сапасына кері әсер етті. Сондықтан инфрақызыл кептіру
үшін оңтайлы температура 60°C және қызанақ ұнтағы үшін 20 сағат кептіру
уақыты, сапаны және қалдық ылғалдылықты 13,1\% сақтайды. Қарбыз ұнтағы
үшін оңтайлы температура 50°C және 37 сағат кептіру уақыты, ылғалдылық
13,4\% сақталады және органолептикалық қасиеттерінің қараюын немесе
жоғалуын көрсетпейді.

Инфрақызыл кептірудің басты артықшылығы - материалдан ылғалды тез және
тиімді түрде кетіру мүмкіндігі, бұл биологиялық белсенді заттармен жұмыс
істегенде маңызды. Инфрақызыл кептіру материалға терең енуді және ішінен
қыздыру арқылы ылғалды тез кетіруді қамтиды. Жылуды беру үшін ауаны
пайдаланатын дәстүрлі әдістерден айырмашылығы, инфрақызыл сәулелену
өнімге тікелей әсер етеді, кептіру процесін жеделдетеді және с дәрумені,
ликопин және басқа антиоксиданттар сияқты белсенді компоненттердің
термиялық ыдырау қаупін азайтады.

Қызанақ және қарбыз ұнтақтарының органолептикалық және физика-химиялық
параметрлерін зерттеу

Қызанақ және қарбыз ұнтақтарының органолептикалық және физика-химиялық
параметрлері (4-кесте) -40°C және --55°C температурада мұздатып кептіру,
сондай-ақ қарбыз ұнтағы үшін 60°C және қызанақ ұнтағы үшін 50°C
температурада инфрақызыл кептіру арқылы зерттелді.

Қарбыз ұнтағының ылғалдылығы мұздатып кептіргеннен кейін 12,5\%, ал
қызанақ ұнтағыныкі 13,3\% болды; инфрақызыл кептіргеннен кейін 13,1\%,
ал қызанақ ұнтағыныкі 13,4\% болды.
\end{multicols}

\tcap{4-кесте. Қарбыз бен қызанақ ұнтақтарының органолептикалық және физика-химиялық қасиеттері}
\begin{longtblr}[
  label = none,
  entry = none,
]{
  width = \linewidth,
  colspec = {Q[198]Q[187]Q[185]Q[175]Q[196]},
  cells = {c},
  cells = {font = \small},
  row{1} = {font=\bfseries\small},
  cell{1}{1} = {r=4}{},
  cell{1}{2} = {c=4}{0.742\linewidth},
  cell{2}{2} = {c=2}{0.372\linewidth},
  cell{2}{4} = {c=2}{0.371\linewidth},
  cell{3}{2} = {c=4}{0.742\linewidth},
  cell{5}{2} = {c=2}{0.372\linewidth},
  cell{5}{4} = {c=2}{0.371\linewidth},
  cell{7}{2} = {c=4}{0.742\linewidth},
  hlines,
  vlines,
}
Көрсеткіш атауы & Зерттелетін ұнтақтар &  &  & \\
 & Арбуз ұнтағы &  & Қызанақ ұнтағы & \\
 & Кептіру әдісі &  &  & \\
 & лиофильді кептіру & Инфрақызыл кептіру & лиофильді кептіру & Инфрақызыл кептіру\\
Дәм және иіс & тәтті, қарбызға тән айқын дәмі және жағымды хош иіс &  & тәтті, қарбызға тән айқын дәмі және жағымды хош иіс & \\
Түс & ашық қызғылт-қызыл & айқын қызыл & ашық кызыл & қанық қызыл\\
Консистенция & ұнақ &  &  & \\
Ылғалдылық, \% & 12,5 & 13,1 & 13,3 & 13,4\\
Титрленетін қышқылдық, \textsuperscript{0} Т & 0,47 & 0,46 & 0,41 & 0,38\\
Белсенді қышқылдық, бірлік & 5,2 & 5,1 & 4,4 & 4,3
\end{longtblr}

\begin{multicols}{2}
Қарбыз ұнтағының мұздатып кептіргеннен және инфрақызыл кептіргеннен
кейін титрленетін қышқылдығы 0,38-0,47 °T аралығында, ал белсенді
қышқылдығы 4,3--5,2 бірлік аралығында болды.

Органолептикалық қасиеттері бойынша қызанақ пен қарбыз ұнтақтары ұнтақ
тәрізді көрінді; Түсі ашық қызылдан ашық қызылға дейін болды; Иісі мен
дәмі кептірілген қызанақ пен қарбызға тән болды, бөгде иістер болмады.

Физикалық-химиялық параметрлер бойынша қарбыз бен қызанақ ұнтақтары
қолайлы шектерде болды. Нәтижелер 4-кестеде келтірілген.

Инфрақызыл кептіру тамақ және фармацевтика өнеркәсібінде көкөністер,
жемістер, дәмдеуіштер және өсімдік сығындылары сияқты кептірілген
өнімдерді өндіру үшін, сондай-ақ дәрі-дәрмектер мен косметика
өндірісінде кеңінен қолданылады. Дегенмен, өнім сапасын сақтау үшін
процестің параметрлерін мұқият бақылау маңызды: инфрақызыл сәулелену
қарқындылығы, температура және кептіру уақыты, себебі шамадан тыс жоғары
температура сезімтал заттарды бұзуы мүмкін.

Қызанақ пен қарбыз ұнтақтарының түсі зерттелді. Барлық үлгілер бір
бастапқы материалдан алынды және таңдалған температура жағдайында
мұздатып кептіру және инфрақызыл кептіру сияқты әртүрлі кептіру әдістері
қолданылды. Түс сипаттамалары CIE Lab* жүйесін пайдалана отырып, Minolta
CM-23d портативті спектрофотометрімен өлшенді, мұндағы L* жарықтық (0 =
қара, 100 = ақ), a* - қызару дәрежесі және b* - сарғыштық дәрежесі.

Қызанақ пен қарбыз ұнтақтары үшін түс стандарттары қол жетімді
болмағандықтан, бұл дақылдардың жаңа піскен жемістері салыстырмалы
зерттеулер үшін бақылау үлгілері ретінде пайдаланылды. Зерттеу барысында
әртүрлі кептіру режимдерінің, соның ішінде ауа ағынының әртүрлі
жылдамдығымен 40, 50 және 60°C температурада инфрақызыл кептірудің,
сондай-ақ әртүрлі температурада мұздатып кептірудің әсері зерттелді.
Барлығы алты үлгі зерттелді:

-үлгі 1 - бақылау (жаңа томаттар);

-үлгі 2 - бақылау (жаңа қарбыз);

-үлгі 3 - инфрақызыл кептіру әдісімен алынған томат ұнтағы;

-үлгі 4 - лиофильді кептіру әдісімен алынған томат ұнтағы;

-үлгі 5 - инфрақызыл кептіру әдісімен алынған қарбыз ұнтағы;

-үлгі 6 - лиофильді кептіру әдісімен алынған қарбыз ұнтағы.

Өлшеулер әр үлгінің он нүктесінде кептіру процесінен кейін бірден
жүргізілді, сонымен қатар түстің қанықтылығы (∆C*) есептелді. Себебі
кептірілген томаттардың түстік стандарты жоқ, бақылау үшін жаңа
томаттардың параметрлері әдебиетте көрсетілген {[}9,10{]}.

Түстің қанықтылығын растау үшін хроматичность немесе CIE LCh* жүйесі
бойынша қанықтылықты есептеу қарастырылуы мүмкін, мұндағы хрома (C*)
мына формуламен есептеледі:

\begin{equation*}
C = \sqrt{a^{*2} + b^{*2}}^*
\end{equation*}

Түстің тұрақтылығы D65 стандартты жарықтандыруында бағаланды - күндізгі
жарықты имитациялау, түс температурасы 6504K, «Ортаңғы солтүстік аспан»
(CIE D65), бейтарап рең.

Қызанақ пен қарбыз ұнтақтарының түс сипаттамаларын салыстырмалы зерттеу
әртүрлі кептіру әдістерін: инфрақызыл кептіру және лиофилизацияны
қолдану арқылы жүргізілді. Түс қанықтылығын бағалау түс параметрлерінің
кептіру әдісі мен процесінің жағдайларына тәуелділігін анықтауға
бағытталған. Нәтижелер қосымшадағы 5-кестеде келтірілген.
\end{multicols}

\tcap{5-кесте. Спектрофотометрде жүргізілген зерттеулердің нәтижелері}
\begin{longtblr}[
  label = none,
  entry = none,
]{
  cells = {c},
  cells = {font = \small},
  row{1} = {font=\bfseries\small},
  cell{1}{1} = {r=2}{},
  cell{1}{2} = {c=4}{},
  hlines,
  vlines,
}
Ұнтақ үлгілері & Спектофотометр көрсеткіштері &  &  & \\
 & \(L\) & \(a\) & \(B\) & түстің қанықтылығы \(С_{ab}\) ед.\\
Үлгі 1 & 38 ,12 & 19.21 & 18,74 & 720,21\\
Үлгі 2 & 34,11 & 18,87 & 16,91 & 642,02\\
Үлгі 3 & 41.17 & 30.19 & 34.03 & 2073,56\\
Үлгі 4 & 50.75 & 22.68 & 21.48 & 975,77\\
Үлгі 5 & 41.42 & 25.22 & 26.35 & 1330,37\\
Үлгі 6 & 58.45 & 20.35 & 21.35 & 869,94
\end{longtblr}

\begin{multicols}{2}
Өлшеу әдісіндегідей, кептіру және температура режимдері режимді
модификациялау сипаттамаларына айтарлықтай әсер етеді: режимде алынған
ұнтақтар (лиофилизация) жоғары жарықтықпен және қызыл түстің
қанықтылығымен (жоғары L* және a* мәндерімен) сипатталады, ал инфрақызыл
сәулеленудің жоғары температурасында кептірілген ұнтақтарда индикаторлар
төмендейді, ал b* компоненті артады, бұл сары-қызғылт сары түске ауысуға
әкеледі.

Лиофильді және инфрақызыл кептіру арқылы алынған қызанақ ұнтақтарындағы
ликопин мөлшерін өзгертуге болады, бұл өнімге әртүрлі әсер етеді.
Зерттеу ең қанық түс инфрақызыл кептіру әдісімен алынған қызанақ
ұнтағында (Cab = 2073,5) және инфрақызыл кептіру әдісімен алынған қарбыз
ұнтағында (Cab = 1330,37) алынғанын анықтады.

Инфрақызыл кептіру мен мұздатып кептірудің үйлесімі жоғары сапалы
қанықтылықты қамтамасыз етеді. Бұл пигмент қосылыстары жақсы сақталған.
Сондай-ақ, ұнтақтарды алудың технологиялық процестері зерттелді және
өнімдегі қанықтылықты сақтау және ликопинді сақтау үшін инфрақызыл
кептіру таңдалды.

Қызанақ және қарбыз ұнтақтарының витаминдік құрамы зерттелді (Кесте 6).
\end{multicols}

\tcap{6 - кесте. Ұнтақтардың (қарбыз және томат) витаминдік құрамы}
\begin{longtblr}[
  label = none,
  entry = none,
]{
  width = 0.8\linewidth,
  colspec = {Q[421]Q[138]Q[138]Q[242]},
  cells = {c},
  cells = {font = \small},
  row{1} = {font=\bfseries\small},
  cell{1}{1} = {r=2}{},
  cell{1}{2} = {c=2}{0.276\linewidth},
  cell{1}{4} = {r=2}{},
  hlines,
  vlines,
}
Көрсеткіштердің атауы, өлшем бірлігі & Ұнтақтар &  & Тәулік қажеттілік (мг)\\
 & қарбыз & қызанақ & \\
Витаминдер, мг/100 г &  &  & \\
В1 & 0,18±0,036 & 0,12±0,02 & 1,1-1,4\\
В2 & 0,07±0,029 & 0,38±0,16 & 1,1-1,6\\
В3 & 4,35±0,87 & 3,84±0,77 & 14-18\\
В5 & 0,32±0,058 & 1,92±0,35 & 5-7\\
В6 & 0,35±0,07 & 0,75±0,15 & 1,2-1,7\\
В9 & 0,010±0,002 & 0,035±0,007 & 0,4\\
С & 84,12±1,09 & 133,05±1,46 & 0,075-0,1\\
А & 1,44 & 88,36 & 0,7-0,9\\
Е & - & 17,13 & 15
\end{longtblr}

\begin{multicols}{2}
Кесте 6 нәтижелеріне сәйкес томат ұнтағының құрамында (мг/100 г): А --
88,36; Е -- 17,13; В1 -- 0,12; В2 -- 0,38; В3 -- 3,84; В5 -- 1,92; В6 --
0,75; В9 -- 0,035; ал қарбыз ұнтағы құрамында (мг/100 г) -- А -- 1,44;
В1 -- 0,18; В2 -- 0,07; В3 -- 4,35; В5 -- 0,32; В6 -- 0,35; В9 -- 0,010.
Ұнтақтардың витаминдік құрамын күнделікті тұтыну бойынша \% есептелді.
Қарбыз ұнтағы күнделікті қажеттілікті қамтамасыз етеді (\%): В1 --
16,4\,\%, В2 -- 6,4\,\%, В3 -- 31,1\,\%, В5 -- 6,4\,\%, В6 -- 2,1\,\%,
В9 -- 2,5\,\%, ал витамин А 100\,\% қамтиды; томат ұнтағы -- В1 --
11\,\%, В2 -- 34,5\,\%, В3 -- 27,4\,\%, В5 -- 38,4\,\%, В6 -- 62,5\,\%,
В9 -- 8,75\,\%, ал витамин А және Е 100\,\% қамтиды. Витаминдер құрамын
талдау және олардың күнделікті қажеттілікті қанағаттандыру қабілетіне
сүйене отырып, мынадай қорытынды жасауға болады: томат және қарбыз
ұнтақтары антиоксиданттық витаминдер (А, Е) және В тобының витаминдері,
әсіресе В6, маңызды болып табылатын метаболизм және нерв жүйесі үшін
жоғары витаминдік құндылығы бар функционалды өнімдер болып табылады.

{\bfseries Қорытынды.} Зерттеу нәтижесінде лиофильді және инфрақызыл
кептірудің оңтайлы режимдері ылғалдылықты, органолептикалық,
физика-химиялық көрсеткіштерді және түсін сақтай отырып, қызанақ пен
қарбыз ұнтағын алуға мүмкіндік беретіні анықталды. Лиофильді кептіру
өнімнің құрылымын, дәмін және биологиялық белсенді заттарын максималды
түрде сақтайды, ал инфрақызыл кептіру түс қанықтылығы жоғары өнімдерді
тез алуға мүмкіндік береді.

Кептіру нәтижелеріне сәйкес, органолептикалық және физика-химиялық
көрсеткіштер бойынша зерттеу нәтижелері келесідей: лиофильді кептіру
режимінде температура --40°C және --55°C болды, ал инфрақызыл кептіру
режимінде қарбыз ұнтағы үшін температура 60°C және қызанақ ұнтағы үшін
50°C тиімді болып таңдалды.

Лиофильді кептіру кезінде қарбыз ұнтағының ылғалдылығы 12,5\%, ал
қызанақ ұнтағының ылғалдылығы 13,3\% болды; инфрақызыл кептіру кезінде
қарбыз ұнтағы 13,1\%, ал қызанақ ұнтағы 13,4\% болды.

Ең қанық түс қызанақ ұнтағын (Cab = 2073.56) және қарбыз ұнтағын (Cab =
1330.37) инфрақызыл кептіру арқылы алынды. Қарбыз ұнтағын мұздатып
кептіру және инфрақызыл кептіру кезінде титрленетін қышқылдық 0.38--0.47
°C аралығында, ал белсенді қышқылдық 4.3--5.2 бірлік аралығында болды.
Органолептикалық көрсеткіштер бойынша, қызанақ және қарбыз ұнтақтары
талапқа сай; түсі ашық қызылдан ашық қызылға дейін; иісі мен дәмі
кептірілген қызанақ пен қарбызға тән, бөгде иістерсіз.

Бұл нәтижелер жоғары сапалы, дәрумендік құндылығы жоғары және
органолептикалық қасиеттері бар функционалды ұнтақ өнімдерін алу үшін
мұздатып кептіру және инфрақызыл кептіру әдістерінің тиімділігін
растайды.

\emph{{\bfseries Қаржыландыру}. Материалдар BR22886613 «Инновациялық
технологияларды дамыту» ғылыми-техникалық бағдарламасы аясында
«Онкологиялық аурулардың алдын алу үшін құны төмен және жоғары сапа
көрсеткіштерімен функционалды тағамдық қоспаларды өндіру технологиясын
жасау» жобасы аясында дайындалған. Ауыл шаруашылығы министрлігінің 101
«Ғылыми зерттеулер мен іс-шараларды бағдарламалық-нысаналы қаржыландыру»
кіші бағдарламасы бойынша 267 «Білім мен ғылыми зерттеулердің
қолжетімділігін арттыру» бюджеттік бағдарламасы бойынша ауыл шаруашылығы
өсімдік шаруашылығы өнімдері мен шикізатын өңдеу және сақтау
технологиялары»; Қазақстан Республикасы 2024-2026 жж.}
\end{multicols}

\begin{center}
{\bfseries Әдебиеттер}
\end{center}

\begin{refs}
1. Ribeiro D., Freitas M., Silva A. M. S., Carvalho F., Fernandes E.
Antioxidant and pro-oxidant activities of carotenoids and their
oxidation products// Food and Chemical Toxicology. -2018. -Vol.120.
-P.681--699. DOI 10.1016/j.fct.2018.07.060.

2. Kulawik A., Cielecka-Piontek J., Zalewski P. The Importance of
Antioxidant Activity for the Health-Promoting Effect of Lycopene//
Nutrients. -2023. -Vol.15(17):3821. DOI 10.3390/nu15173821.

3. Bin-Jumah M. N., Nadeem M. S., Gilani S. J., Mubeen B., Ullah I.,
Alzarea S. I., Ghoneim M. M., Alshehri S., Al-Abbasi F. A., Kazmi I.
Lycopene: A Natural Arsenal in the War against Oxidative Stress and
Cardiovascular Diseases// Antioxidants. -2022. -Vol.11(2): 232. DOI
10.3390/antiox11020232.

4. Saini, R. K., Rengasamy, K. R. R., Mahomoodally, F. M., \& Keum, Y.-S.
Protective effects of lycopene in cancer, cardiovascular, and
neurodegenerative diseases: An update on epidemiological and
mechanistic perspectives// Pharmacological Research. -2020. -Vol.155:
104730. DOI 10.1016/j.phrs.2020.104730.

5. Antonuccio P., Micali A., Puzzolo D., Romeo C., Vermiglio G.,
Squadrito V., Freni J., Pallio G., Trichilo V., Righi M., Irrera N.,
Altavilla D., Squadrito F., MariniH. R., Minutoli L. Nutraceutical
Effects of Lycopene in Experimental Varicocele: An ``In Vivo'' Model
to Study Male Infertility// Nutrients. -2020. -Vol.12(5): 1536. DOI
10.3390/nu12051536.

6. Ni Y., Zhuge F., Nagashimada M., Nagata N., Xu L., Yamamoto S., Fuke
N., Ushida Y., Suganuma H., Kaneko S., Ota T. Lycopene prevents the
progression of lipotoxicity-induced nonalcoholic steatohepatitis by
decreasing oxidative stress in mice// Free Radical Biology and
Medicine. -2020. -Vol.152. -P.571-582. DOI
10.1016/j.freeradbiomed.2019.11.036.

7. Eslami E., Carpentieri S., Pataro G., Ferrari G. A Comprehensive
Overview of Tomato Processing By-Product Valorization by Conventional
Methods versus Emerging Technologies//Foods. -2022. -Vol.12(1): 166.
DOI 10.3390/foods12010166.

8. Lado J., Zacarias J., Rodrigo M. J., Zacarías L. Visualization of
Carotenoid-Storage Structures in Fruits by Transmission Electron
Microscopy// Methods in Molecular Biology. -2020. -Vol.2083. -P.
235-244.
\href{https://doi.org/10.1007/978-1-4939-9952-1_18}{DOI:10.1007/978-1-4939-9952-1\_18}.

9. Cucu T., Huvaere K., Van Den Bergh M.A., et al. A simple and fast HPLC
method to determine lycopene in foods// Food Analytical Methods.
-2012. -Vol.5. -P.1221-1228.
\href{https://doi.org/10.1007/s12161-011-9354-6}{DOI
10.1007/s12161-011-9354-6}.

10. Khalid M., Saeed R., Bilal M., Iqbal H. M. N., Huang D. Biosynthesis
and biomedical perspectives of carotenoids with special reference to
human health-related applications// Biocatalysis and Agricultural
Biotechnology. -2019. -Vol.17. -P.399-407.
\href{https://doi.org/10.1016/j.bcab.2018.11.027}{DOI
10.1016/j.bcab.2018.11.027}.
\end{refs}

\begin{info}
\emph{{\bfseries Авторлар туралы мәліметтер}}

Жумалиева Г.Е. -- техника ғылымдарының кандидаты, «Қазақ қайта өңдеу
және тағам өнеркәсіптері ғылыми-зерттеу институты» ЖШС Алматы қаласы
филиалы, Алматы, Казақстан, e-mail: guljan\_7171@mail.ru;

Чоманов У.Ч. -- т.ғ.д., «Қазақ қайта өңдеу және тағам өнеркәсіптері
ғылыми-зерттеу институты» ЖШС Алматы қаласы филиалы, Алматы, Казақста,
e-mail: сhomanov@mail.ru;

Байзақова А.Б. - жаратылыстану ғылымдарының магистрі, ғылыми қызметкер
«Қазақ қайта өңдеу және тағам өнеркәсіптері ғылыми-зерттеу институты»
ЖШС Алматы қаласы филиалы, Алматы, Казақстан, e-mail:
ainel0026@gmail.com;

Шоман А.К.- техника ғылымдарының магистрі, аға ғылыми қызметкер,
«Қазақ қайта өңдеу және тағам өнеркәсіптері ғылыми-зерттеу институты»
ЖШС Алматы қаласы филиалы, Алматы, Казақстан, e-mail:
shoman\_a@mail.ru.

\emph{{\bfseries Information about the authors}}

Zhumalieva G.E. - candidate of technical sciences, Almaty Branch of
LLP ``Kazakh Research Institute of Processing and Food Industry'',
Almaty, Kazakhstan, e-mail: guljan\_7171@mail.ru;

Chomanov U.Ch. - doctor of technical sciences, Almaty Branch of LLP
``Kazakh Research Institute of Processing and Food Industry'', Almaty,
Kazakhstan, e-mail: сhomanov@mail.ru;

Baizakova A.B. - master of natural sciences, research fellow, Almaty
Branch of LLP ``Kazakh Research Institute of Processing and Food
Industry'', Almaty, Kazakhstan, e-mail: ainel0026@gmail.com;

Shoman A.K. - master of technical sciences, Almaty Branch of LLP
``Kazakh Research Institute of Processing and Food Industry'', Almaty,
Kazakhstan, e-mail: shoman\_a@mail.ru.
\end{info}
